\documentclass[11pt]{article}

\usepackage[T1]{fontenc}
\usepackage[utf8]{inputenc}
\usepackage{lmodern}
\usepackage{amsmath,amssymb,amsfonts,amsthm}
\usepackage[margin=1in]{geometry}
\usepackage{microtype}
\usepackage{xcolor}
\usepackage[colorlinks=true,linkcolor=teal,urlcolor=teal]{hyperref}

\newcommand{\R}{\mathbb R}
\newcommand{\C}{\mathbb C}
\newcommand{\Hc}{\mathsf H}
\newcommand{\G}{\mathsf G}
\newcommand{\F}{\mathcal F_{\mathrm{orb}}}
\newcommand{\T}{\mathcal T_1}
\newcommand{\K}{\mathcal K}
\newcommand{\A}{\mathfrak A}
\newcommand{\At}{\widetilde{\mathfrak A}}
\newcommand{\M}{\mathfrak M}
\newcommand{\Tin}{\Theta^{\mathrm{in}}}
\newcommand{\Ttar}{\Theta^{\mathrm{tar}}}
\newcommand{\Uin}{\mathcal U^{\mathrm{in}}}
\newcommand{\Uout}{\mathcal U^{\mathrm{out}}}
\newcommand{\Tr}{\operatorname{Tr}}
\newcommand{\dd}{\,\mathrm d}
\newcommand{\B}{\mathrm B}

\newtheorem*{lemmaC}{Lemma C.2 (rephrased)}

\title{A Short Guide to Lemma C.2\\
\large F\o lner averaging, compactness, and exact covariance}
\author{}
\date{}

\begin{document}
\maketitle

\begin{abstract}
Lemma C.2 turns an optimization over arbitrary hybrid channels, averaged on
larger and larger parameter cubes, into an optimization at the origin over
exactly covariant maps.  The proof combines four ideas: group averaging,
the F\o lner property of Euclidean cubes, weak-* compactness of channel
adjoints, and upper semicontinuity of fidelity.  This note develops those
ideas and explains why the limiting map may lose trace when mass escapes in
the noncompact classical coordinate.
\end{abstract}

\section{What the lemma accomplishes}

The hybrid Gaussian experiment has a classical coordinate
$h\in\Hc$ and a quantum displacement coordinate $z\in\C^{r_d}$.  Write
\[
  \G=\Hc\times\C^{r_d},\qquad \xi=(h,z).
\]
The input and desired target form two covariant families
\[
  \Tin_\xi=\Uin_\xi(\Tin_0),
  \qquad
  \Ttar_\xi=\Uout_\xi(\Ttar_0).
\]
For a channel $\Lambda$, the paper studies the average payoff
\[
  \frac1{|\Omega_L|}\int_{\Omega_L}
  \B\bigl(\Lambda(\Tin_\xi),\Ttar_\xi\bigr)\dd\xi
\]
over a large centered cube $\Omega_L\subset\G$.  An arbitrary channel can
behave differently at different parameter values.  Lemma C.2 says that its
asymptotic average performance cannot exceed the performance at the origin
of some map $\Gamma$ satisfying the exact intertwining relation
\[
  \Gamma\Uin_\xi=\Uout_\xi\Gamma
  \qquad(\xi\in\G).
\]
Thus a noncompact, many-parameter optimization is reduced to a single-state
optimization with a strong symmetry constraint.

\section{The hybrid state spaces}

Let $\F$ be the multimode Fock space.  An ordinary hybrid state is a
Bochner-integrable field of trace-class operators,
\[
  \mathfrak T=L^1(\Hc;\T(\F)).
\]
Its natural test algebra is
\[
  \A=C_0(\Hc;\K(\F)),
  \qquad \M=\A^*.
\]
Positive elements of $\M$ are finite positive $\T(\F)$-valued Radon
measures.  Allowing outputs in $\M$, rather than only in $\mathfrak T$, is
important: weak limits may have singular classical parts such as point
masses.

The group actions used in the paper are
\begin{align*}
  (\Uin_\xi R)(x)
  &=D(z)R(x-h)D(z)^*,\\
  (\Uout_\xi\Xi)(E)
  &=D(\sqrt\gamma z)\,
    \Xi(E-\sqrt\gamma h)\,
    D(\sqrt\gamma z)^*.
\end{align*}
They translate the classical component and conjugate the quantum component
by Weyl displacements.  Both preserve positivity and total trace.  Root
fidelity is invariant under these reversible actions.

Since $\A$ is nonunital, the proof also uses its unitization
\[
  \At=\A\oplus\C I.
\]
A normalized hybrid measure has a canonical state extension to $\At$.  The
extra quotient character of $\At$ represents a point at infinity.  This
compactifies the state space sufficiently for weak-* extraction.

\section{Background idea I: F\o lner sets}

A locally compact group is amenable if it possesses sets whose relative
boundary becomes negligible under every fixed translation.  Here the group
is a Euclidean vector group, and centered cubes form a F\o lner family:
\begin{equation}\label{eq:folner}
  \frac{|(\Omega_L+v)\mathbin\triangle\Omega_L|}{|\Omega_L|}
  \longrightarrow0
  \qquad(L\to\infty)
\end{equation}
for every fixed $v\in\G$.  The symbol $\triangle$ denotes symmetric
difference.

Geometrically, translation by a fixed vector changes only slabs near the
faces of a large cube.  Their volume is of order $L^{\dim\G-1}$, whereas the
cube has volume of order $L^{\dim\G}$.  Their ratio therefore tends to zero.
Equation \eqref{eq:folner} is the mechanism that turns approximate covariance
of averaged channels into exact covariance of a limit.

\section{Background idea II: compactness through adjoints}

There is no useful norm compactness for a sequence of infinite-dimensional
channels.  Instead one passes to adjoints.  A trace-preserving completely
positive map induces a unital completely positive contraction
\[
  \widehat\Lambda^*: \At\longrightarrow\mathfrak T^*,
  \qquad
  \mathfrak T^*=L^\infty(\Hc;\mathcal B(\F)).
\]
The unit balls of the dual space $\mathfrak T^*$ are weak-* compact.  Because
$\At$ and the predual $\mathfrak T$ are separable, weak-* compactness is
metrizable on bounded subsets.  One may therefore choose a countable dense
operator system in $\At$ and diagonalize to obtain pointwise weak-* convergence
of the adjoints.

Complete positivity survives this limit.  Indeed, for every matrix size
$r$, positivity of a map means
\[
  [A_{kl}]\geq0
  \quad\Longrightarrow\quad
  [\widehat\Lambda^*(A_{kl})]\geq0.
\]
The positive cone of the dual matrix space is weak-* closed, so this property
passes to the limiting adjoint at every $r$.  This is why the proof checks
matrix positivity, rather than only ordinary positivity.

\section{Statement of Lemma C.2}

\begin{lemmaC}
Let $\Omega_L$ be the centered cubes in $\G$.  Then
\begin{equation}\label{eq:statement}
  \limsup_{L\to\infty}\sup_\Lambda
  \frac1{|\Omega_L|}\int_{\Omega_L}
  \B\bigl(\Lambda(\Tin_\xi),\Ttar_\xi\bigr)\dd\xi
  \leq
  \sup_\Gamma \B\bigl(\Gamma(\Tin_0),\Ttar_0\bigr),
\end{equation}
where $\Lambda$ ranges over hybrid channels and $\Gamma$ ranges over
completely positive, trace-nonincreasing maps
$\mathfrak T\to\M$ such that
\[
  \Gamma\Uin_\xi=\Uout_\xi\Gamma
  \qquad(\xi\in\G).
\]
The output of $\Gamma$ is allowed to be a measure rather than an ordinary
$L^1$ density.
\end{lemmaC}

The direction of \eqref{eq:statement} matters.  The lemma supplies an upper
bound: every nearly optimal sequence of arbitrary channels produces an
exactly covariant limiting competitor whose origin fidelity is at least the
limiting average payoff.

\section{Proof, step by step}

\subsection{Step 1: average the channel in a moving frame}

Define
\begin{equation}\label{eq:average-map}
  \overline\Lambda_L
  =\frac1{|\Omega_L|}\int_{\Omega_L}
  \Uout_{-\xi}\circ\Lambda\circ\Uin_\xi\dd\xi.
\end{equation}
For a parameter $\xi$, the input is first shifted from the origin to $\xi$,
then processed by $\Lambda$, and finally shifted back in the output
coordinates.  Thus every parameter point is viewed from the origin before
the results are averaged.

Because $\Tin_\xi=\Uin_\xi(\Tin_0)$ and
$\Ttar_\xi=\Uout_\xi(\Ttar_0)$, invariance of root fidelity gives
\[
  \B\bigl(
    \Uout_{-\xi}\Lambda(\Tin_\xi),\Ttar_0
  \bigr)
  =\B\bigl(\Lambda(\Tin_\xi),\Ttar_\xi\bigr).
\]
Root fidelity is jointly concave, hence concave in its first argument with
the target fixed.  Applying this to \eqref{eq:average-map} yields
\begin{equation}\label{eq:concavity}
  \frac1{|\Omega_L|}\int_{\Omega_L}
  \B\bigl(\Lambda(\Tin_\xi),\Ttar_\xi\bigr)\dd\xi
  \leq
  \B\bigl(\overline\Lambda_L(\Tin_0),\Ttar_0\bigr).
\end{equation}
This is the first decisive reduction: an average of fidelities is bounded by
one fidelity at the origin.

\subsection{Step 2: extract a weak-* limiting map}

Choose $L_j\to\infty$ and channels whose left side in
\eqref{eq:statement} approaches the limsup.  For brevity, use
$\overline\Lambda_j$ for their averages.  The unitized adjoints
\[
  \widehat\Lambda_j^*: \At\to\mathfrak T^*
\]
are unital completely positive contractions.  The diagonal compactness
argument described above gives, after passing to a subsequence,
\[
  \widehat\Lambda_j^*(A)
  \stackrel{w^*}{\longrightarrow}
  \widehat\Gamma^*(A)
  \qquad(A\in\At).
\]
Contractivity first defines the limit on a countable dense operator system
and then extends it uniquely to all of $\At$.  Weak-* closedness of every
matrix positive cone makes $\widehat\Gamma^*$ completely positive, and
convergence at $I$ makes it unital.

Restrict $\widehat\Gamma^*$ to the ideal $\A$ and define
$\Gamma:\mathfrak T\to\M$ by
\begin{equation}\label{eq:duality}
  \langle\Gamma(R),A\rangle
  =\langle R,\widehat\Gamma^*(A)\rangle,
  \qquad A\in\A.
\end{equation}
This restriction is the measure-valued limiting map needed on the right side
of \eqref{eq:statement}.

\subsection{Step 3: explain why only trace nonincrease remains}

Let $(e_\ell)$ be an increasing contractive approximate identity for $\A$.
For $R\geq0$, the total mass of the positive measure $\Gamma(R)$ is
\[
  \Tr\Gamma(R)
  =\sup_\ell\langle\Gamma(R),e_\ell\rangle.
\]
Using \eqref{eq:duality}, positivity, and $0\leq e_\ell\leq I$ gives
\begin{align*}
  \Tr\Gamma(R)
  &=\sup_\ell\langle R,\widehat\Gamma^*(e_\ell)\rangle\\
  &\leq\langle R,\widehat\Gamma^*(I)\rangle
   =\Tr R.
\end{align*}
Thus $\Gamma$ is trace nonincreasing.

The apparent loss is not a defect in the argument.  The unitized limit is
still normalized, but part of its classical mass may sit at the quotient
point at infinity.  Restricting from $\At$ to $\A$ discards precisely that
part.  Requiring the restricted limit to remain trace preserving would
incorrectly assume tightness of all output measures.

\subsection{Step 4: the limiting map is exactly covariant}

For a fixed $v\in\G$, translating the integration variable in
\eqref{eq:average-map} shows that the two maps
$\overline\Lambda_j\Uin_v$ and
$\Uout_v\overline\Lambda_j$ differ only on the symmetric difference of the
two averaging cubes.  Consequently, for $R\in\mathfrak T$,
\begin{equation}\label{eq:cov-error}
  \left\|
  \overline\Lambda_j(\Uin_vR)
  -\Uout_v\overline\Lambda_j(R)
  \right\|_{\M}
  \leq
  \frac{|(\Omega_{L_j}+v)\triangle\Omega_{L_j}|}{|\Omega_{L_j}|}
  \|R\|_1.
\end{equation}
The F\o lner ratio tends to zero by \eqref{eq:folner}.  Pairing
\eqref{eq:cov-error} with $A\in\A$ and passing to the pointwise weak-* limit
gives
\[
  \Gamma\Uin_vR=\Uout_v\Gamma(R).
\]
Since $v$ was arbitrary, $\Gamma$ is exactly covariant under all of $\G$.

\subsection{Step 5: pass the fidelity payoff to the limit}

Weak-* convergence does not generally imply convergence of fidelity.  The
needed one-sided statement follows from the Alberti--Uhlmann variational
formula on a unital $C^*$-algebra:
\begin{equation}\label{eq:AU}
  \B(\phi,\psi)
  =\inf_{A>0\ \mathrm{invertible}}
  \frac12\bigl(\phi(A)+\psi(A^{-1})\bigr).
\end{equation}
For fixed $\psi$, each expression inside the infimum is weak-* continuous in
$\phi$.  An infimum of continuous functions is upper semicontinuous, so
\[
  \phi_j\stackrel{w^*}{\longrightarrow}\phi
  \quad\Longrightarrow\quad
  \limsup_j\B(\phi_j,\psi)\leq\B(\phi,\psi).
\]
Apply this on $\At$ to the unitized origin outputs.  Lemma C.1 then removes
any mass at infinity: the target is normalized on $\A$, so its overlap with
the quotient component is zero.  Therefore
\[
  \limsup_j
  \B\bigl(\overline\Lambda_j(\Tin_0),\Ttar_0\bigr)
  \leq
  \B\bigl(\Gamma(\Tin_0),\Ttar_0\bigr).
\]
Combining this inequality with \eqref{eq:concavity} proves
\eqref{eq:statement}.

\section{A finite-dimensional analogy}

For a finite group $G$, one may average a channel exactly:
\[
  \overline\Lambda
  =\frac1{|G|}\sum_{g\in G}
  \Uout_{g^{-1}}\Lambda\Uin_g.
\]
This average is already exactly covariant, so no limit or compactness
argument is required.  Lemma C.2 is the noncompact analogue.  A large
F\o lner cube plays the role of an approximate uniform distribution on
$\G$, and weak-* compactness supplies a limiting average because no finite
translation-invariant probability measure exists on $\G$.

\section{Dependencies and role in the paper}

Lemma C.2 uses:
\begin{itemize}
  \item joint concavity and invariance of root fidelity;
  \item the F\o lner property of Euclidean cubes;
  \item weak-* compactness and separability to extract channel adjoints;
  \item weak-* closedness of matrix positive cones;
  \item the Alberti--Uhlmann variational formula; and
  \item Lemma C.1 to discard fidelity-invisible mass at infinity.
\end{itemize}

Its output is the starting point for Lemma C.3 and Proposition C.4.  Lemma
C.3 describes the structure of an exactly covariant hybrid map through a
residual subinstrument.  Proposition C.4 then bounds the origin fidelity of
every such map by
\[
  B_{\mathrm{univ}}(\gamma,p)
  =B_{\mathrm{cl}}(\gamma,d)B_{\mathrm{orb}}(\gamma,p).
\]
Together, C.2 and C.4 give the upper bound in the hybrid Gaussian flat-prior
theorem.

\section{Minimal prerequisite checklist}

The conceptual core requires only covariance, concavity of root fidelity,
and the fact that the boundary-to-volume ratio of a large cube vanishes.
For the rigorous compactness step, one additionally needs dual Banach spaces,
the weak-* topology, adjoints of completely positive maps, unitization of a
nonunital $C^*$-algebra, and the variational characterization
\eqref{eq:AU}.  The main message is simple: average a channel in the moving
frame, take a compact limit, and use vanishing boundary effects to make the
limit exactly symmetric.

\end{document}
